%% kapa: exam=1 round=37 year=2026 period=1 form=A subject=real_estate subject_ko=부동산학원론 seq=23 number=103
%% source: Q-Net 감정평가사 기출문제 (cst00309) artlSeq=5258924 file=raw/1008/5258924/2026년 제37회 감정평가사 1차 1교시.pdf
%% conversion: kordoc 4.24.51 -> tools/gen1.py (auto); review=proofread-sonnet; flags=none
\begin{question}{103}
위험회피형 투자자가 다음 조건의 두 자산 X, Y를 대상으로 포트폴리오를 구성하려고 한다. 이에 관한 설명으로 옳은 것은? (단, 주어진 조건에 한함)
\begin{framedtext}
○ X와 Y의 기대수익률의 표준편차: $\sigma_{x}<\sigma_{y}$ \par
○ X와 Y의 투자비중: $w_{x}\geq 0,\ w_{y}\geq 0,\ w_{x}+w_{y}=1$ \par
\end{framedtext}
\begin{choices}
  \item 최소분산포트폴리오에서 X의 투자비중은 0.5보다 크다.
  \item X와 Y의 수익률의 상관계수가 작을수록, 포트폴리오의 기대수익률은 커진다.
  \item X와 Y의 수익률의 상관계수가 ‘0’이면, 포트폴리오의 분산은 ‘0’이 된다.
  \item X와 Y의 수익률의 상관계수가 ‘0’보다 크면, 포트폴리오의 표준편차는 X의 표준편차보다 크거나 같다.
  \item X와 Y의 수익률의 상관계수가 $-1$이면, 무위험포트폴리오를 만들기 위한 자산의 투자비중은 $w_{x}=\dfrac{\sigma_{x}}{\sigma_{x}+\sigma_{y}}$, $w_{y}=\dfrac{\sigma_{y}}{\sigma_{x}+\sigma_{y}}$이다.
\end{choices}
\end{question}
\answer{1}
